\documentclass[11pt]{article}
% ===========================================================================
%  Shared preamble for all MPM2D worksheets — input right after \documentclass.
%  (Files beginning with "_" are skipped by mpm2d-worksheet-publish.mjs.)
% ===========================================================================
\usepackage[letterpaper,top=2.2cm,bottom=1.9cm,left=1.6cm,right=1.6cm,headheight=28pt,footskip=24pt]{geometry}
\usepackage{amsmath,amssymb}
\usepackage{cancel}
\usepackage[T1]{fontenc}
\usepackage[utf8]{inputenc}
\usepackage{lmodern}
\usepackage{xcolor}
\usepackage[most]{tcolorbox}
\usepackage{enumitem}
\usepackage{multicol}
\usepackage{fancyhdr}
\usepackage{tikz}
\usetikzlibrary{calc}
\usepackage{pgfplots}
\pgfplotsset{compat=1.18}

\definecolor{exblue}{HTML}{4A90E2}
\definecolor{exbg}{HTML}{E6F3FF}
\definecolor{qorange}{HTML}{E69138}
\definecolor{qbg}{HTML}{FFF7CC}
\definecolor{keygray}{HTML}{475569}
\definecolor{keybg}{HTML}{F1F5F9}
\definecolor{iagreen}{HTML}{14653B}
\definecolor{titleblue}{HTML}{1D4ED8}
% Rotating palette so each worked example gets its own colour.
\definecolor{ex0}{HTML}{2563A0}\definecolor{ex1}{HTML}{3B7D3B}\definecolor{ex2}{HTML}{8A5A00}
\definecolor{ex3}{HTML}{A3327A}\definecolor{ex4}{HTML}{6D28A3}\definecolor{ex5}{HTML}{B08900}
\definecolor{ex6}{HTML}{0E7490}\definecolor{ex7}{HTML}{9A3412}\definecolor{ex8}{HTML}{15803D}

\pagestyle{fancy}
\fancyhf{}
\fancyhead[L]{\footnotesize NAME: \underline{\hspace{3.4cm}}}
\fancyhead[C]{\textbf{Integration Academy}\\[-2pt]{\scriptsize Math Simplified \textperiodcentered{} Success Amplified}}
\fancyhead[R]{\footnotesize MPM2D}
\fancyfoot[L]{\scriptsize dr.merie@IntegrationAcademy.ca}
\fancyfoot[C]{\scriptsize\thepage}
\fancyfoot[R]{\scriptsize IntegrationAcademy.ca}
\renewcommand{\headrulewidth}{1.2pt}
\renewcommand{\footrulewidth}{0.4pt}
\setlength{\parindent}{0pt}
\setlength{\parskip}{4pt}

% Examples are NOT breakable, so each example + its solution stays on one page.
\newtcolorbox{example}[2]{colframe=#1,colback=#1!7,coltitle=white,colbacktitle=#1,fonttitle=\bfseries,title={#2},arc=3pt,boxrule=1pt}
\newtcolorbox{qbox}[1]{colframe=qorange,colback=qbg,coltitle=white,colbacktitle=qorange,fonttitle=\bfseries,title={#1},arc=3pt,breakable,boxrule=1pt}
\newtcolorbox{keybox}{colframe=keygray,colback=keybg,coltitle=white,colbacktitle=keygray,fonttitle=\bfseries,title={$\checkmark$\ Answer Key},arc=3pt,breakable,boxrule=1pt}
\newtcolorbox{recapbox}{colframe=keygray,colback=keybg,coltitle=white,colbacktitle=keygray,fonttitle=\bfseries,title={Question Recap},arc=3pt,breakable,boxrule=1pt}
\newtcolorbox{ideabox}{colframe=titleblue,colback=titleblue!6,coltitle=white,colbacktitle=titleblue,fonttitle=\bfseries,title={The Big Ideas},arc=3pt,breakable,boxrule=1.2pt}
\newtcolorbox{learnbox}{colframe=iagreen,colback=iagreen!5,coltitle=white,colbacktitle=iagreen,fonttitle=\bfseries,title={The Big Picture},arc=3pt,breakable,boxrule=1.2pt}
\newcommand{\lh}[1]{\par\smallskip{\bfseries\color{iagreen}#1}\par\nobreak\smallskip}

\newcommand{\soln}{\par\smallskip\textit{\textbf{Solution.}}\ }
\newcommand{\workspace}[1]{\par\vspace{3pt}\fcolorbox{black!30}{white}{\begin{minipage}[t][#1][t]{\dimexpr\linewidth-2\fboxsep\relax}\ \end{minipage}}\par}
\newcommand{\sech}[1]{\par\vspace{8pt}{\Large\bfseries\color{iagreen}#1}\par\nobreak\vspace{2pt}{\color{black!20}\hrule height1pt}\vspace{6pt}}

% A compact coordinate plot sized for an example box: \eplot{xmin}{xmax}{ymin}{ymax}{body}
\newcommand{\eplot}[5]{\begin{center}\begin{tikzpicture}
\begin{axis}[width=6.8cm,height=4.4cm,axis lines=middle,xlabel={\tiny$x$},ylabel={\tiny$y$},
grid=both,grid style={gray!16},xmin=#1,xmax=#2,ymin=#3,ymax=#4,
xtick distance=2,ytick distance=2,tick label style={font=\tiny},axis line style={-},samples=120]
#5
\end{axis}\end{tikzpicture}\end{center}}

% A sine-curve plot (degrees on x, 0..360): \splot{ymin}{ymax}{body}
\newcommand{\splot}[3]{\begin{center}\begin{tikzpicture}
\begin{axis}[width=8.6cm,height=4.6cm,axis lines=middle,xlabel={\tiny$x\,(^\circ)$},ylabel={\tiny$y$},
grid=both,grid style={gray!16},xmin=0,xmax=360,ymin=#1,ymax=#2,
xtick distance=90,ytick distance=1,tick label style={font=\tiny},axis line style={-},samples=180]
#3
\end{axis}\end{tikzpicture}\end{center}}

% A small coordinate plot: \plot{xmin}{xmax}{ymin}{ymax}{body}
\newcommand{\plot}[5]{\begin{center}\begin{tikzpicture}
\begin{axis}[width=8.4cm,height=6cm,axis lines=middle,xlabel={\scriptsize$x$},ylabel={\scriptsize$y$},
grid=both,grid style={gray!16},xmin=#1,xmax=#2,ymin=#3,ymax=#4,
xtick distance=2,ytick distance=2,tick label style={font=\tiny},axis line style={-}]
#5
\end{axis}\end{tikzpicture}\end{center}}

% Right triangle (angle theta at bottom-left, right angle at bottom-right):
%   \rtri{drawBase}{drawHeight}{oppLabel}{adjLabel}{hypLabel}
\newcommand{\rtri}[5]{\begin{center}\begin{tikzpicture}[scale=0.85]
\coordinate (A) at (0,0);\coordinate (B) at (#1,0);\coordinate (C) at (#1,#2);
\draw[thick,exblue] (A)--(B)--(C)--cycle;
\draw[black] ($(B)+(-0.32,0)$)--($(B)+(-0.32,0.32)$)--($(B)+(0,0.32)$);
\node[below] at ($(A)!0.5!(B)$){#4};
\node[right] at ($(B)!0.5!(C)$){#3};
\node[above left] at ($(A)!0.5!(C)$){#5};
\node at (0.7,0.22){$\theta$};
\end{tikzpicture}\end{center}}

% General (oblique) triangle with vertex labels and opposite-side labels:
%   \gtri{Alabel}{Blabel}{Clabel}{aSide}{bSide}{cSide}
\newcommand{\gtri}[6]{\begin{center}\begin{tikzpicture}[scale=0.85]
\coordinate (A) at (0,0);\coordinate (B) at (5,0);\coordinate (C) at (1.6,3);
\draw[thick,exblue] (A)--(B)--(C)--cycle;
\node[below left] at (A){#1};\node[below right] at (B){#2};\node[above] at (C){#3};
\node[below] at ($(A)!0.5!(B)$){#6};
\node[right] at ($(B)!0.5!(C)$){#4};
\node[left] at ($(A)!0.5!(C)$){#5};
\end{tikzpicture}\end{center}}

\fancyhead[R]{\footnotesize MCV4U \textperiodcentered{} 2.3}
\begin{document}

\begin{center}
{\LARGE\bfseries\color{titleblue}Worksheet 2.3 --- The Chain Rule}\\[2pt]
{\small Grade 12 Calculus and Vectors (MCV4U) \textperiodcentered{} Unit 2: Derivative Rules}
\end{center}

\begin{ideabox}
$\dfrac{d}{dx}f(g(x))=f'(g(x))\cdot g'(x)$ — derivative of the outside (inside untouched) times the derivative of the inside.
\begin{itemize}[leftmargin=*,itemsep=2pt]
  \item Chain rule: $f'(g(x))\cdot g'(x)$.
  \item Power of a function: $\dfrac{d}{dx}[g(x)]^n=n[g(x)]^{n-1}g'(x)$.
  \item Don't forget the inner derivative $g'(x)$.
\end{itemize}
\end{ideabox}

\begin{learnbox}
\lh{Differentiating a composition}The \textbf{chain rule} handles a function inside a function: $\dfrac{d}{dx}f(g(x))=f'(g(x))\cdot g'(x)$.
\lh{Outside times inside}Differentiate the outer function while keeping the inner intact, then multiply by the derivative of the inner.
\lh{Layer by layer}For nested functions, apply the chain rule repeatedly, working from the outermost layer inward.
\end{learnbox}

\sech{Worked Examples}

\begin{example}{ex0}{Example 1 --- Power of a function}
Differentiate $f(x)=(x^2+1)^3$.\soln
\textbf{Step 1 (identify outer and inner):} outer is $(\,\cdot\,)^3$, inner is $g(x)=x^2+1$ (so $g'(x)=2x$).\\[3pt]
\textbf{Step 2 (differentiate the outer, leaving the inner alone):} $3(x^2+1)^2$.\\[3pt]
\textbf{Step 3 (multiply by the inner derivative):} $3(x^2+1)^2\cdot2x$.\\[3pt]
\textbf{Conclusion:} $f'(x)=6x(x^2+1)^2$.
\end{example}

\begin{example}{ex1}{Example 2 --- Radical}
Differentiate $f(x)=\sqrt{x^2+1}$.\soln
\textbf{Step 1 (rewrite as a power):} $f(x)=(x^2+1)^{1/2}$, with inner $g(x)=x^2+1$ (so $g'(x)=2x$).\\[3pt]
\textbf{Step 2 (differentiate the outer power):} $\tfrac12(x^2+1)^{-1/2}$.\\[3pt]
\textbf{Step 3 (multiply by the inner derivative):} $\tfrac12(x^2+1)^{-1/2}\cdot2x$.\\[3pt]
\textbf{Step 4 (simplify):} the $2$'s cancel, and $(x^2+1)^{-1/2}=\dfrac1{\sqrt{x^2+1}}$.\\[3pt]
\textbf{Conclusion:} $f'(x)=\dfrac{x}{\sqrt{x^2+1}}$.
\end{example}

\begin{example}{ex2}{Example 3 --- Linear inside}
Differentiate $f(x)=(3x-2)^5$.\soln
\textbf{Step 1 (identify outer and inner):} outer is $(\,\cdot\,)^5$, inner is $g(x)=3x-2$ (so $g'(x)=3$).\\[3pt]
\textbf{Step 2 (differentiate the outer):} $5(3x-2)^4$.\\[3pt]
\textbf{Step 3 (multiply by the inner derivative):} $5(3x-2)^4\cdot3$.\\[3pt]
\textbf{Conclusion:} $f'(x)=15(3x-2)^4$.
\end{example}

\begin{example}{ex3}{Example 4 --- Cubic inside}
Differentiate $f(x)=(x^3+x)^4$.\soln
\textbf{Step 1 (identify outer and inner):} outer is $(\,\cdot\,)^4$, inner is $g(x)=x^3+x$ (so $g'(x)=3x^2+1$).\\[3pt]
\textbf{Step 2 (differentiate the outer):} $4(x^3+x)^3$.\\[3pt]
\textbf{Step 3 (multiply by the inner derivative):} $4(x^3+x)^3(3x^2+1)$.\\[3pt]
\textbf{Conclusion:} $f'(x)=4(x^3+x)^3(3x^2+1)$.
\end{example}

\begin{example}{ex4}{Example 5 --- Chain inside a product}
Differentiate $f(x)=x(2x+1)^3$.\soln
\textbf{Step 1 (identify the two factors for the product rule):} $u=x$ (so $u'=1$) and $v=(2x+1)^3$.\\[3pt]
\textbf{Step 2 (differentiate $v$ with the chain rule):} outer $3(\,\cdot\,)^2$, inner derivative $2$, so $v'=3(2x+1)^2\cdot2=6(2x+1)^2$.\\[3pt]
\textbf{Step 3 (apply the product rule $u'v+uv'$):} $(1)(2x+1)^3+x\cdot6(2x+1)^2$.\\[3pt]
\textbf{Step 4 (factor out the common $(2x+1)^2$):} $(2x+1)^2\left[(2x+1)+6x\right]=(2x+1)^2(8x+1)$.\\[3pt]
\textbf{Conclusion:} $f'(x)=(2x+1)^2(8x+1)$.
\end{example}

\begin{example}{ex5}{Example 6 --- Power of a function}
Differentiate $f(x)=(x^2-4)^3$.\soln
\textbf{Step 1 (identify outer and inner):} outer is $(\,\cdot\,)^3$, inner is $g(x)=x^2-4$ (so $g'(x)=2x$).\\[3pt]
\textbf{Step 2 (differentiate the outer, multiply by the inner derivative):} $3(x^2-4)^2\cdot2x$.\\[3pt]
\textbf{Conclusion:} $f'(x)=6x(x^2-4)^2$.
\end{example}

\begin{example}{ex6}{Example 7 --- Linear inside}
Differentiate $f(x)=(5x+1)^4$.\soln
\textbf{Step 1 (identify outer and inner):} outer is $(\,\cdot\,)^4$, inner is $g(x)=5x+1$ (so $g'(x)=5$).\\[3pt]
\textbf{Step 2 (differentiate the outer, multiply by the inner derivative):} $4(5x+1)^3\cdot5$.\\[3pt]
\textbf{Conclusion:} $f'(x)=20(5x+1)^3$.
\end{example}

\begin{example}{ex7}{Example 8 --- Radical}
Differentiate $f(x)=\sqrt{2x+1}$.\soln
\textbf{Step 1 (rewrite as a power):} $f(x)=(2x+1)^{1/2}$, with inner derivative $2$.\\[3pt]
\textbf{Step 2 (differentiate the outer, multiply by the inner derivative):} $\tfrac12(2x+1)^{-1/2}\cdot2$.\\[3pt]
\textbf{Step 3 (simplify, the $2$'s cancel):} $(2x+1)^{-1/2}$.\\[3pt]
\textbf{Conclusion:} $f'(x)=\dfrac1{\sqrt{2x+1}}$.
\end{example}

\begin{example}{ex8}{Example 9 --- Reciprocal power}
Differentiate $f(x)=\dfrac1{(x+1)^2}$.\soln
\textbf{Step 1 (rewrite as a power):} $f(x)=(x+1)^{-2}$, with inner derivative $1$.\\[3pt]
\textbf{Step 2 (differentiate the outer, multiply by the inner derivative):} $-2(x+1)^{-3}\cdot1$.\\[3pt]
\textbf{Step 3 (rewrite back as a fraction):} $-2(x+1)^{-3}=\dfrac{-2}{(x+1)^3}$.\\[3pt]
\textbf{Conclusion:} $f'(x)=\dfrac{-2}{(x+1)^3}$.
\end{example}

\clearpage
\sech{Your Turn --- Show Your Work}
For each question, show all steps clearly in the space provided.

\begin{qbox}{Question 1}
Differentiate $f(x)=(x^2+3)^4$.\workspace{2.4cm}
\end{qbox}

\begin{qbox}{Question 2}
Differentiate $f(x)=(4x-1)^3$.\workspace{2.4cm}
\end{qbox}

\begin{qbox}{Question 3}
Differentiate $f(x)=\sqrt{x^2+9}$.\workspace{2.4cm}
\end{qbox}

\begin{qbox}{Question 4}
Differentiate $f(x)=(x^3-x)^5$.\workspace{2.4cm}
\end{qbox}

\begin{qbox}{Question 5}
Differentiate $f(x)=\dfrac1{(2x+1)^2}$.\workspace{2.4cm}
\end{qbox}

\begin{qbox}{Question 6}
Differentiate $f(x)=(1-x^2)^4$.\workspace{2.4cm}
\end{qbox}

\begin{qbox}{Question 7}
Differentiate $f(x)=\sqrt{3x-2}$.\workspace{2.4cm}
\end{qbox}

\begin{qbox}{Question 8}
Differentiate $f(x)=(x^2+x+1)^3$.\workspace{2.4cm}
\end{qbox}

\begin{qbox}{Question 9}
Differentiate $f(x)=x^2(x+1)^3$.\workspace{2.4cm}
\end{qbox}

\begin{qbox}{Question 10}
Differentiate $f(x)=(2x+5)^6$.\workspace{2.4cm}
\end{qbox}

\begin{qbox}{Question 11}
Differentiate $f(x)=\dfrac{1}{\sqrt{x+4}}$.\workspace{2.4cm}
\end{qbox}

\begin{qbox}{Question 12}
Differentiate $f(x)=(x^2-1)^{1/2}$.\workspace{2.4cm}
\end{qbox}

\begin{qbox}{Question 13 --- Challenge}
Differentiate $f(x)=\dfrac{(x+1)^2}{x}$ (chain + quotient).\workspace{3cm}
\end{qbox}

\clearpage
\sech{Question Recap \& Answer Key}

\begin{recapbox}
\begin{multicols}{2}
\small
\begin{enumerate}[leftmargin=*,itemsep=3pt]
  \item Differentiate $f(x)=(x^2+3)^4$.
  \item Differentiate $f(x)=(4x-1)^3$.
  \item Differentiate $f(x)=\sqrt{x^2+9}$.
  \item Differentiate $f(x)=(x^3-x)^5$.
  \item Differentiate $f(x)=\dfrac1{(2x+1)^2}$.
  \item Differentiate $f(x)=(1-x^2)^4$.
  \item Differentiate $f(x)=\sqrt{3x-2}$.
  \item Differentiate $f(x)=(x^2+x+1)^3$.
  \item Differentiate $f(x)=x^2(x+1)^3$.
  \item Differentiate $f(x)=(2x+5)^6$.
  \item Differentiate $f(x)=\dfrac{1}{\sqrt{x+4}}$.
  \item Differentiate $f(x)=(x^2-1)^{1/2}$.
  \item Differentiate $f(x)=\dfrac{(x+1)^2}{x}$ (chain + quotient).
\end{enumerate}
\end{multicols}
\end{recapbox}

\begin{keybox}
\begin{multicols}{2}
\small
\begin{enumerate}[leftmargin=*,itemsep=3pt]
  \item $8x(x^2+3)^3$
  \item $12(4x-1)^2$
  \item $\dfrac{x}{\sqrt{x^2+9}}$
  \item $5(x^3-x)^4(3x^2-1)$
  \item $\dfrac{-4}{(2x+1)^3}$
  \item $-8x(1-x^2)^3$
  \item $\dfrac3{2\sqrt{3x-2}}$
  \item $3(x^2+x+1)^2(2x+1)$
  \item $2x(x+1)^3+3x^2(x+1)^2$
  \item $12(2x+5)^5$
  \item $-\tfrac12(x+4)^{-3/2}$
  \item $\dfrac{x}{\sqrt{x^2-1}}$
  \item $\dfrac{2x(x+1)\cdot x-(x+1)^2}{x^2}=\dfrac{(x+1)(x-1)}{x^2}$
\end{enumerate}
\end{multicols}
\end{keybox}

\end{document}
